% gauge_from_cd3_highschool_ja.tex
% Compile with: xelatex gauge_from_cd3_highschool_ja.tex
\documentclass[12pt,a4paper]{article}

\usepackage{xeCJK}
\setCJKmainfont{Hiragino Mincho ProN}
\setCJKsansfont{Hiragino Kaku Gothic ProN}

\usepackage{amsmath,amssymb}
\usepackage{geometry}
\usepackage{hyperref}
\usepackage{booktabs}
\usepackage{graphicx}
\usepackage{xcolor}
\usepackage[most]{tcolorbox}

\geometry{margin=2.5cm}

% Custom tcolorbox styles
\newtcolorbox{keyidea}[1][]{
  colback=blue!5!white,
  colframe=blue!60!black,
  fonttitle=\bfseries\large,
  title={#1},
  arc=3mm,
  boxrule=1pt
}

\newtcolorbox{analogy}[1][]{
  colback=green!5!white,
  colframe=green!50!black,
  fonttitle=\bfseries,
  title={#1},
  arc=2mm,
  boxrule=0.8pt
}

\newtcolorbox{mathbox}[1][]{
  colback=orange!5!white,
  colframe=orange!60!black,
  fonttitle=\bfseries,
  title={#1},
  arc=2mm,
  boxrule=0.8pt
}

\newtcolorbox{honestbox}[1][]{
  colback=red!5!white,
  colframe=red!50!black,
  fonttitle=\bfseries,
  title={#1},
  arc=2mm,
  boxrule=0.8pt
}

\newtcolorbox{wowbox}[1][]{
  colback=purple!5!white,
  colframe=purple!50!black,
  fonttitle=\bfseries\large,
  title={#1},
  arc=3mm,
  boxrule=1.2pt
}

\newcommand{\Spec}{\operatorname{Spec}}

\title{\textbf{\Huge たった一つの数が宇宙を支配する}\\[8pt]
\Large cd = 3 から時空も力もダークエネルギーも}

\author{Wright Brothers Collaboration\\
{\normalsize --- 高校生のための解説ガイド ---}}
\date{2026年4月}

\begin{document}
\maketitle

\begin{tcolorbox}[colback=yellow!10!white, colframe=yellow!60!black,
  fonttitle=\bfseries\Large, title={この論文のメッセージ}, arc=4mm]
宇宙の最も基本的な性質---空間が3次元であること、
力が3種類あること、ダークエネルギーの量---が、
たった一つの数学的な数「$\mathrm{cd} = 3$」から導かれるかもしれない。
\end{tcolorbox}

\vspace{1em}

\tableofcontents

\newpage

%======================================================================
\section{一つの数がすべてを支配する：cd = 3}
%======================================================================

\begin{keyidea}[cd = 3 とは何か？]
$\mathrm{cd}(\Spec(\mathbb{Z})) = 3$

これは「整数の世界の次元は3である」という意味の数学的な定理です。
1960年代にArtin（アルティン）とVerdier（ヴェルディエ）という
数学者が証明しました。

$\Spec(\mathbb{Z})$は「すべての素数が住んでいる空間」、
cdは「コホモロジー次元」（その空間の"本当の次元"を測る道具）です。
\end{keyidea}

\begin{analogy}[DNAのアナロジー]
DNAにはたった4つの塩基（A, T, G, C）しかありません。
でもこの4文字で、すべての生命の設計図が書かれています。

同じように、$\mathrm{cd} = 3$というたった一つの数が、
宇宙の「設計図」の大部分を決めているかもしれません：

\begin{center}
\begin{tabular}{ll}
\toprule
DNAの4塩基 & $\to$ すべての生命 \\
$\mathrm{cd} = 3$ & $\to$ 時空の次元 + 自然界の力 + ダークエネルギー \\
\bottomrule
\end{tabular}
\end{center}
\end{analogy}

\subsection{なぜ「3」なのか？}

$\mathrm{cd} = 3$は人間が選んだ数ではありません。
数学的に証明された事実です。
素数（2, 3, 5, 7, 11, \ldots）の分布の深い性質から来ています。

もし素数の分布が違っていたら、cdの値も違ったかもしれません。
でも実際の数学では、$\mathrm{cd}(\Spec(\mathbb{Z})) = 3$は
変えようがない定理です。

%======================================================================
\section{cd = 3 から時空へ（簡単な復習）}
%======================================================================

\begin{keyidea}[空間3次元 + 時間1次元 = 4次元時空]
\begin{center}
\textbf{空間の次元} = cd = 3（上下、左右、前後）

\textbf{時間の次元} = 1（量子重力の境界条件から）

\textbf{時空の次元} = 3 + 1 = \textbf{4}
\end{center}
\end{keyidea}

これは前の論文で詳しく説明しました。
ポイントは：
\begin{itemize}
  \item 空間が3次元なのは、$\Spec(\mathbb{Z})$が
    「3次元の空間」のように振る舞うから。
  \item 時間が1次元なのは、Wheeler--DeWitt方程式という
    量子重力の基礎方程式が、空間の「外側」に
    ちょうど1つの方向を追加するから。
\end{itemize}

\begin{analogy}[部屋のアナロジー]
3次元の部屋（cd = 3）の中にいると想像してください。
あなたは壁・天井・床（3方向）を見ることができます。
そして時間は、部屋の「外」を流れています。
部屋の次元（3）が時空の次元（3+1=4）を決めるのです。
\end{analogy}

%======================================================================
\section{cd = 3 から力へ（今回の新発見！）}
%======================================================================

ここからが今回の論文の一番大事なところです。

\subsection{自然界の3つの力}

物理学によれば、自然界には（重力を除いて）
3つの基本的な力があります：

\begin{center}
\begin{tabular}{llll}
\toprule
力 & 対称性 & 日常での例 & 担い手 \\
\midrule
\textbf{電磁気力} & $\mathrm{U}(1)$ & 光、電気、磁石 & 光子（1個） \\
\textbf{弱い力} & $\mathrm{SU}(2)$ & 放射性崩壊 & $W^+, W^-, Z$（3個） \\
\textbf{強い力} & $\mathrm{SU}(3)$ & 原子核を結びつける & グルーオン（8個） \\
\bottomrule
\end{tabular}
\end{center}

なぜこの3つなのか？ なぜ4つ目や5つ目がないのか？
これは物理学の大きな謎でした。

\subsection{「3次元の家」のアナロジー}

\begin{analogy}[3部屋の家]
$\Spec(\mathbb{Z})$は「3次元の家」だと思ってください。

この家には1次元、2次元、3次元の「部屋」があります：
\begin{itemize}
  \item \textbf{1次元の部屋}（線のような部屋）$\to$ ここに電磁気力が住む
  \item \textbf{2次元の部屋}（平面のような部屋）$\to$ ここに弱い力が住む
  \item \textbf{3次元の部屋}（立体の部屋）$\to$ ここに強い力が住む
\end{itemize}

4次元の部屋は？ \textbf{存在できません}。
家が3次元だから、4次元以上の部屋は物理的に入りません。
だから4つ目の力はないのです！
\end{analogy}

\subsection{数学的に言うと}

\begin{mathbox}[Langlands対応：数論と物理の架け橋]
数学には「Langlands（ラングランズ）対応」という深い理論があります。
この理論によると：

\begin{center}
\begin{tabular}{ccc}
\toprule
数論の世界 & $\longleftrightarrow$ & 物理の世界 \\
\midrule
$\mathrm{GL}(1)$の対象 & $\longleftrightarrow$ & $\mathrm{U}(1)$（電磁気力） \\
$\mathrm{GL}(2)$の対象 & $\longleftrightarrow$ & $\mathrm{SU}(2)$（弱い力） \\
$\mathrm{GL}(3)$の対象 & $\longleftrightarrow$ & $\mathrm{SU}(3)$（強い力） \\
$\mathrm{GL}(4)$の対象 & $\longleftrightarrow$ & \textbf{存在しない！}（cd=3だから） \\
\bottomrule
\end{tabular}
\end{center}

$\mathrm{GL}(n)$は「$n$次元の行列の群」です。
$n$が$\mathrm{cd} = 3$を超えると、
$\Spec(\mathbb{Z})$の中に「収まらない」のです。
\end{mathbox}

\begin{wowbox}[もしcd = 4だったら？]
もし$\Spec(\mathbb{Z})$の次元が4だったら：
\begin{itemize}
  \item 空間は4次元（4方向に動ける！）
  \item 力は4種類（$\mathrm{SU}(4)$という追加の力がある！）
  \item 追加の力には15個のゲージボソンが必要
\end{itemize}
そんな宇宙は私たちの宇宙とはまったく違うものになります。
$\mathrm{cd} = 3$であることが、私たちの宇宙を
私たちの宇宙たらしめているのです。
\end{wowbox}

%======================================================================
\section{12個のゲージボソン = $2/|B_2|$}
%======================================================================

\begin{keyidea}[力の担い手の数]
3つの力を合わせると、ゲージボソン（力を伝える粒子）は
全部で何個でしょうか？

\begin{center}
\textbf{光子} 1個 + \textbf{弱ボソン} 3個 + \textbf{グルーオン} 8個 = \textbf{12個}
\end{center}
\end{keyidea}

この「12」は偶然の数ではありません。
数学には「Bernoulli数」という有名な数列があり、
その2番目の値は$B_2 = 1/6$です。すると：
\[
  \frac{2}{|B_2|} = \frac{2}{1/6} = 12.
\]

\begin{mathbox}[Riemannゼータ関数との繋がり]
Riemannゼータ関数$\zeta(s)$を$s = -1$で評価すると：
\[
  \zeta(-1) = -\frac{1}{12}
\]
よって：
\[
  \text{ゲージボソンの数} = \frac{1}{|\zeta(-1)|} = 12
\]

ゼータ関数は素数の分布を記述する関数です。
ゲージボソンの数が素数の性質と繋がっているのです！

しかもこの$1/12$は、ダークエネルギーの公式にも現れます。
偶然でしょうか？ 私たちはそうは思いません。
\end{mathbox}

%======================================================================
\section{ダークエネルギーの3つの材料}
%======================================================================

\begin{keyidea}[ダークエネルギーの公式]
\[
  \Omega_\Lambda = \underbrace{\frac{4}{3}}_{\text{重力}} \times
                   \underbrace{2\pi}_{\text{幾何学}} \times
                   \underbrace{\frac{1}{12}}_{\text{算術}} =
                   \frac{2\pi}{9} \approx 0.698
\]

観測値：$\Omega_\Lambda^{\text{obs}} \approx 0.689$（Planck衛星、2018年）

一致度：約$1.3\%$の差！
\end{keyidea}

3つの因子はそれぞれ別の起源を持ちます：

\begin{analogy}[ケーキのアナロジー]
ダークエネルギーはケーキを焼くようなものです。
3つの材料がすべて必要です：

\begin{center}
\begin{tabular}{lll}
\toprule
ケーキの材料 & ダークエネルギーの因子 & 値 \\
\midrule
\textbf{小麦粉}（基礎） & 算術的因子（ゼータ関数） & $1/12$ \\
\textbf{水}（形を作る） & 幾何学的因子（$\pi$） & $2\pi$ \\
\textbf{熱}（焼き上げる） & 重力的因子（Friedmann方程式） & $4/3$ \\
\bottomrule
\end{tabular}
\end{center}

小麦粉だけではケーキにならない。
水だけでもダメ。熱だけでもダメ。
3つが組み合わさってはじめてダークエネルギーが生まれます。
\end{analogy}

\subsection{各因子の意味}

\begin{tcolorbox}[colback=gray!5!white, colframe=gray!50!black, title={因子1：重力 = $4/3$}]
アインシュタインの重力理論によると、宇宙の膨張は
Friedmann方程式で記述されます。
4次元時空では、この方程式に$4/3$という係数が現れます。

もし宇宙が5次元だったら$5/4$、6次元だったら$6/5$になります。
$4/3$であることは、まさに時空が4次元であることの反映です。
\end{tcolorbox}

\begin{tcolorbox}[colback=gray!5!white, colframe=gray!50!black, title={因子2：幾何学 = $2\pi$}]
$\pi = 3.14159\ldots$は円周率です。
この因子は、Riemannゼータ関数の「ガンマ因子」から来ます：
\[
  \xi_\infty(-1) = \pi^{1/2}\,\Gamma(-1/2) = -2\pi
\]
$\pi$は実数の世界に固有の数です。
後で見るように、これがダークエネルギーの存在に
決定的な役割を果たします。
\end{tcolorbox}

\begin{tcolorbox}[colback=gray!5!white, colframe=gray!50!black, title={因子3：算術 = $1/12$}]
$\zeta(-1) = -1/12$はRiemannゼータ関数の値です。
ここでは$p = 2$のEuler因子を調整した値を使います：
\[
  \zeta_{\neg 2}(-1) = \frac{1}{12}
\]
これは素数の分布に由来する純粋に算術的な量です。
ゲージボソンの数（12個）の逆数でもあります。
\end{tcolorbox}

%======================================================================
\section{なぜ$\pi$が大事か：実数の魔法}
%======================================================================

\begin{keyidea}[実数の世界 vs $p$進数の世界]
数学では、有理数$\mathbb{Q}$（分数の世界）を
「完備化」する方法が2種類あります：

\begin{center}
\begin{tabular}{lll}
\toprule
完備化 & 結果 & 特徴 \\
\midrule
ふつうの距離 & 実数 $\mathbb{R}$ & 滑らか、連続、$\pi$がある \\
素数$p$の距離 & $p$進数 $\mathbb{Q}_p$ & フラクタル的、不連続、$\pi$がない \\
\bottomrule
\end{tabular}
\end{center}
\end{keyidea}

\subsection{$\pi$は実数にしかない}

$\pi$は円の周長と直径の比です。
円は「滑らかな曲線」であり、実数$\mathbb{R}$の世界でしか
意味を持ちません。

$p$進数の世界では、空間はフラクタル的で、
「滑らかな円」が描けません。
したがって$\pi$という概念自体が存在しないのです。

\subsection{$\pi$がないとダークエネルギーもない}

\begin{wowbox}[驚くべき帰結]
ダークエネルギーの公式：
\[
  \Omega_\Lambda = \frac{4}{3} \times \underbrace{2\pi}_{\text{この因子！}} \times \frac{1}{12}
\]

もし$\pi$がなかったら？ この因子がゼロになり、
$\Omega_\Lambda = 0$（ダークエネルギーなし）。

つまり：
\begin{center}
\textbf{「私たちの宇宙にダークエネルギーがあるのは、\\
実数が存在するからである。」}
\end{center}

仮想的な「$p$進的宇宙」では、
ダークエネルギーはゼロ---宇宙は加速膨張しません。
\end{wowbox}

\begin{analogy}[音楽のアナロジー]
実数の世界は「連続的な音」（バイオリンのように滑らかに
音程を変えられる）。

$p$進数の世界は「離散的な音」（ピアノのように決まった音しか出せない、
しかもフラクタル的に配置されている）。

ダークエネルギーは「滑らかな振動」から生まれるので、
バイオリン（実数）の世界にはあるけど、
ピアノ（$p$進数）の世界にはないのです。
\end{analogy}

%======================================================================
\section{真空 vs 物質 = Eisenstein vs カスプ}
%======================================================================

\subsection{保型形式の2つの種類}

数論には「保型形式」という美しい関数のクラスがあります。
保型形式は2つの種類に分かれます：

\begin{center}
\begin{tabular}{lll}
\toprule
種類 & 数学的性質 & 物理的対応 \\
\midrule
\textbf{Eisenstein級数} & 滑らか、ゆっくり変化 & 真空（何もない空間） \\
\textbf{カスプ形式} & 振動する、波のよう & 物質（粒子） \\
\bottomrule
\end{tabular}
\end{center}

\begin{analogy}[海のアナロジー]
\begin{itemize}
  \item \textbf{Eisenstein級数} = 海面の平均的な高さ（潮位）。
    ゆっくり変化し、どこでもだいたい同じ。
    これは宇宙の「真空エネルギー」（ダークエネルギー）に対応します。
  \item \textbf{カスプ形式} = 海面の波。
    上下に振動し、エネルギーを運ぶ。
    これは宇宙の「物質」（粒子）に対応します。
\end{itemize}

海面 = Eisenstein（真空）+ 波（物質）、
ちょうど宇宙 = 真空エネルギー + 物質 と同じです！
\end{analogy}

\subsection{Langlands対応でつながる}

カスプ形式には「Galois表現」という数論的な対象が付随します。
これはLanglands対応を通じて、物質粒子の性質を符号化します。

一方、Eisenstein級数にはGalois表現がありません（技術的には「可約」です）。
これは真空が「粒子でない」ことに対応します。

\begin{mathbox}[Wright Brothersの枠組み]
Wright Brothers（ライト兄弟）の飛行機をイメージしてください：

\begin{center}
\begin{tabular}{ll}
\textbf{左翼} & = Eisenstein = 真空 = ダークエネルギー（非摂動論的） \\
\textbf{右翼} & = カスプ = 物質 = 粒子（摂動論的、Langlands対応） \\
\textbf{本体} & = $\Spec(\mathbb{Z})$（すべてを支える算術的構造） \\
\end{tabular}
\end{center}

両翼があってはじめて飛べる---
真空と物質の両方があってはじめて宇宙は「飛ぶ」のです。
\end{mathbox}

%======================================================================
\section{正直な限界}
%======================================================================

\begin{honestbox}[科学は正直でなければならない]
この理論にはいくつかの重要な限界があります。
正直に認めることが科学の基本です。

\begin{itemize}
  \item \textbf{「$n \le \mathrm{cd}$」はヒューリスティック（経験則）であり、
    厳密に証明された定理ではない。}
    $n$次元のGalois表現がcd次元を超えられないという主張は、
    数学的に完全には証明されていません。
    もっともらしい根拠はありますが、厳密な証明はまだです。

  \item \textbf{$\mathrm{GL}(3)$のLanglands対応は未完成。}
    $\mathrm{GL}(2)$（弱い力に対応）のLanglands対応は
    Wilesらにより確立されていますが、
    $\mathrm{GL}(3)$（強い力に対応）は研究途上です。

  \item \textbf{導手（レベル）= 質量 は機能しない。}
    カスプ形式の「導手」が粒子の質量に対応するという
    自然な予想は、現実的な値を与えません。

  \item \textbf{$1.3\%$のずれ。}
    予言値$0.698$と観測値$0.689$には約$1.3\%$のずれがあります。
    これは許容範囲内ですが、何かが足りない可能性もあります。

  \item \textbf{力学がない。}
    ゲージ群（どんな力があるか）は決まりますが、
    ラグランジアン（力がどう働くか）は導出されていません。
\end{itemize}
\end{honestbox}

%======================================================================
\section{全体像：$\Spec(\mathbb{Z})$からすべてへ}
%======================================================================

\begin{wowbox}[cd = 3 から導かれるもの]
\begin{center}
\Large
\textbf{cd = 3}
\end{center}

\begin{center}
$\Downarrow$
\end{center}

\begin{center}
\begin{tabular}{rl}
\textbf{時空次元：} & $d = \mathrm{cd} + 1 = 3 + 1 = 4$ \\[6pt]
\textbf{ゲージ群：} & $\mathrm{U}(1) \times \mathrm{SU}(2) \times \mathrm{SU}(3)$
  （$n \le 3$） \\[6pt]
\textbf{ゲージボソン：} & $1 + 3 + 8 = 12 = 2/|B_2|$ \\[6pt]
\textbf{ダークエネルギー：} & $\displaystyle\Omega_\Lambda
  = \frac{4}{3}\times 2\pi \times \frac{1}{12} = \frac{2\pi}{9}$ \\
\end{tabular}
\end{center}
\end{wowbox}

\subsection{宇宙のソースコードは短いかもしれない}

プログラマーは「美しいコードは短い」と言います。
もしこの理論が正しければ、宇宙の「ソースコード」は
驚くほど短いことになります：

\begin{tcolorbox}[colback=black!90!white, colframe=black,
  colupper=green!70!white]
\texttt{\# 宇宙のソースコード（疑似コード）}\\[4pt]
\texttt{cd = etale\_cohomological\_dimension(Spec(Z))}\\
\texttt{assert cd == 3 ~~~\# これは定理}\\[4pt]
\texttt{spacetime\_dim = cd + 1 ~~~\# = 4}\\
\texttt{gauge\_group = product(SU(n) for n in 1..cd)}\\
\texttt{dark\_energy = d/(d-1) * 2*pi * 1/12}\\[4pt]
\texttt{run\_universe(spacetime\_dim, gauge\_group, dark\_energy)}
\end{tcolorbox}

もちろん、これは大幅に簡略化しています。
でも本質的なメッセージは：

\begin{center}
\textbf{「宇宙の複雑さの根源は、$\Spec(\mathbb{Z})$という\\
極めてシンプルな数学的対象にあるかもしれない。」}
\end{center}

\subsection{まだ分からないこと}

この理論が答えていない大きな問いがたくさんあります：
\begin{itemize}
  \item なぜ電子の質量はあの値なのか？
  \item なぜ世代（電子、ミューオン、タウ）が3つあるのか？
  \item 重力はどう組み込まれるのか？
  \item なぜ物質が反物質より多いのか？
\end{itemize}

これらは今後の研究課題です。
でも、出発点が見つかったかもしれない---
それが$\mathrm{cd}(\Spec(\mathbb{Z})) = 3$なのです。

\subsection{高校生のあなたへ}

\begin{tcolorbox}[colback=blue!3!white, colframe=blue!40!black,
  fonttitle=\bfseries, title={数学と物理学の未来}]
この論文で紹介した内容は、数論（素数の理論）と物理学
（宇宙の理論）が深い所でつながっているかもしれない、
という大きな予感を伝えるものです。

もしあなたがこの話に興味を持ったなら、
以下の分野を勉強してみてください：
\begin{itemize}
  \item \textbf{代数的整数論}：素数やゼータ関数の世界
  \item \textbf{代数幾何学}：$\Spec(\mathbb{Z})$の住む世界
  \item \textbf{表現論}：$\mathrm{GL}(n)$や$\mathrm{SU}(n)$の世界
  \item \textbf{素粒子物理学}：標準模型の世界
  \item \textbf{宇宙論}：ダークエネルギーの世界
\end{itemize}

これらの分野が一つに融合する日が来るかもしれません。
その日を作るのは、もしかするとあなたかもしれません。
\end{tcolorbox}

%======================================================================
\section*{用語集}
%======================================================================

\begin{center}
\begin{tabular}{lp{10cm}}
\toprule
用語 & 意味 \\
\midrule
$\Spec(\mathbb{Z})$ & すべての素数が「点」として住んでいる幾何学的な空間 \\
cd（コホモロジー次元） & 空間の「本当の次元」を測る数学的道具 \\
エタールコホモロジー & 代数的な空間の「形」を調べる道具 \\
ゲージボソン & 力を伝える粒子（光子、$W$ボソン、$Z$ボソン、グルーオン） \\
$\mathrm{GL}(n)$ & $n\times n$の正則行列全体の群 \\
$\mathrm{SU}(n)$ & $n\times n$のユニタリ行列（行列式1）の群 \\
Langlands対応 & $\mathrm{GL}(n)$の保型表現とGalois表現を結ぶ予想 \\
Bernoulli数 & 数論に現れる重要な有理数の列。$B_2 = 1/6$ \\
$\zeta(-1)$ & Riemannゼータ関数の$s=-1$での値。$= -1/12$ \\
ダークエネルギー & 宇宙の加速膨張を引き起こす謎のエネルギー \\
$\Omega_\Lambda$ & ダークエネルギーが宇宙全体のエネルギーに占める割合 \\
保型形式 & ある種の対称性を持つ特別な関数 \\
Eisenstein級数 & 滑らかに変化する保型形式（真空に対応） \\
カスプ形式 & 振動する保型形式（物質に対応） \\
$p$進数 & 素数$p$に基づく「別の距離」で完備化した数の体系 \\
Arakelov幾何学 & 実数と$p$進数を対等に扱う幾何学 \\
\bottomrule
\end{tabular}
\end{center}

\end{document}
